%%%%%%%%%%%%%%%%%%%%%%%%%%%%%%%%%%%%%%%%%%%%%%%%%%%%%%%%%%%%
% 5.16 APPROXIMATION THEORY
% The Composition of Advanced Mathematics
%%%%%%%%%%%%%%%%%%%%%%%%%%%%%%%%%%%%%%%%%%%%%%%%%%%%%%%%%%%%

\section{Approximation Theory}
\label{sec:approximation-theory}

\prerequisite{
Real analysis, linear algebra, calculus, functional analysis, Fourier
analysis, and basic numerical methods.
}

\learningobjectives{
By the end of this section, the reader should be able to:
\begin{itemize}
    \item explain the fundamental goals of approximation theory;
    \item distinguish interpolation from approximation;
    \item understand approximation in different norms;
    \item apply the Weierstrass approximation theorem;
    \item understand the Stone--Weierstrass theorem;
    \item work with polynomial approximation;
    \item construct Bernstein polynomial approximations;
    \item understand best approximation in normed spaces;
    \item use orthogonal projection for least-squares approximation;
    \item work with orthogonal polynomials;
    \item distinguish uniform, least-squares, and minimax approximation;
    \item understand Chebyshev polynomials and equioscillation;
    \item recognize approximation rates and smoothness relationships;
    \item understand Jackson-type approximation estimates;
    \item work with piecewise polynomial and spline approximation;
    \item understand rational approximation and Pad\'e approximation;
    \item connect Fourier approximation with general approximation theory;
    \item recognize approximation by positive linear operators;
    \item understand density of approximation spaces;
    \item connect approximation theory with numerical analysis,
          optimization, data fitting, and scientific computing.
\end{itemize}
}

%%%%%%%%%%%%%%%%%%%%%%%%%%%%%%%%%%%%%%%%%%%%%%%%%%%%%%%%%%%%
\subsection{What Is Approximation Theory?}
\label{subsec:what-is-approximation-theory}
%%%%%%%%%%%%%%%%%%%%%%%%%%%%%%%%%%%%%%%%%%%%%%%%%%%%%%%%%%%%

Approximation theory studies how complicated mathematical objects can be
replaced by simpler ones while controlling the resulting error.

A typical problem begins with a function \(f\) and an approximation
family

\[
\mathcal{A}.
\]

The goal is to find

\[
g\in\mathcal{A}
\]

such that

\[
\boxed{
f\approx g.
}
\]

The approximation family might consist of

\[
\boxed{
\text{polynomials},
\quad
\text{trigonometric polynomials},
\quad
\text{splines},
\quad
\text{rational functions},
\quad
\text{finite-dimensional subspaces}.
}
\]

The quality of the approximation depends on how error is measured.

%%%%%%%%%%%%%%%%%%%%%%%%%%%%%%%%%%%%%%%%%%%%%%%%%%%%%%%%%%%%
\subsection{Approximation Error}
\label{subsec:approximation-error}
%%%%%%%%%%%%%%%%%%%%%%%%%%%%%%%%%%%%%%%%%%%%%%%%%%%%%%%%%%%%

Given a norm \(\|\cdot\|\), the approximation error is

\[
\boxed{
\|f-g\|.
}
\]

Different norms measure different aspects of the error.

For continuous functions on a compact interval, the uniform norm is

\[
\boxed{
\|f\|_\infty
=
\max_{x\in[a,b]}|f(x)|.
}
\]

The corresponding approximation error is

\[
\boxed{
\|f-g\|_\infty
=
\max_{x\in[a,b]}|f(x)-g(x)|.
}
\]

%%%%%%%%%%%%%%%%%%%%%%%%%%%%%%%%%%%%%%%%%%%%%%%%%%%%%%%%%%%%
\subsection{Approximation in \(L^p\)}
\label{subsec:approximation-lp}
%%%%%%%%%%%%%%%%%%%%%%%%%%%%%%%%%%%%%%%%%%%%%%%%%%%%%%%%%%%%

For

\[
1\leq p<\infty,
\]

the \(L^p\) error is

\[
\boxed{
\|f-g\|_p
=
\left(
\int_\Omega
|f(x)-g(x)|^p\,dx
\right)^{1/p}.
}
\]

In particular,

\[
\boxed{
\|f-g\|_2^2
=
\int_\Omega
|f-g|^2\,dx
}
\]

measures mean-square error.

Approximation in \(L^2\) is especially important because \(L^2\) is a
Hilbert space.

%%%%%%%%%%%%%%%%%%%%%%%%%%%%%%%%%%%%%%%%%%%%%%%%%%%%%%%%%%%%
\subsection{Interpolation Versus Approximation}
\label{subsec:interpolation-versus-approximation}
%%%%%%%%%%%%%%%%%%%%%%%%%%%%%%%%%%%%%%%%%%%%%%%%%%%%%%%%%%%%

Interpolation requires exact agreement at selected points:

\[
\boxed{
p(x_j)=f(x_j).
}
\]

Approximation generally does not require exact agreement anywhere.

Instead, it seeks to make an overall error small:

\[
\boxed{
\|f-p\|
\text{ small}.
}
\]

Thus

\[
\boxed{
\text{interpolation}
=
\text{exact at specified nodes},
}
\]

while

\[
\boxed{
\text{approximation}
=
\text{globally close according to a criterion}.
}
\]

%%%%%%%%%%%%%%%%%%%%%%%%%%%%%%%%%%%%%%%%%%%%%%%%%%%%%%%%%%%%
\subsection{Best Approximation}
\label{subsec:best-approximation}
%%%%%%%%%%%%%%%%%%%%%%%%%%%%%%%%%%%%%%%%%%%%%%%%%%%%%%%%%%%%

Let \(X\) be a normed vector space and let \(M\subseteq X\).

\begin{definitionbox}{Best Approximation}
An element \(g^\ast\in M\) is a best approximation to \(f\in X\) from
\(M\) if

\[
\boxed{
\|f-g^\ast\|
=
\inf_{g\in M}\|f-g\|.
}
\]
\end{definitionbox}

The quantity

\[
\boxed{
E(f,M)
=
\inf_{g\in M}\|f-g\|
}
\]

is the distance from \(f\) to the approximation set \(M\).

%%%%%%%%%%%%%%%%%%%%%%%%%%%%%%%%%%%%%%%%%%%%%%%%%%%%%%%%%%%%
\subsection{Existence and Uniqueness}
\label{subsec:best-approximation-existence}
%%%%%%%%%%%%%%%%%%%%%%%%%%%%%%%%%%%%%%%%%%%%%%%%%%%%%%%%%%%%

A best approximation does not automatically exist for every subset of
every normed space.

However, important existence results hold when the approximation set has
additional structure.

For example, in a Hilbert space, every nonempty closed convex set has a
unique nearest point.

In particular, if \(M\) is a closed linear subspace of a Hilbert space,
then every \(f\) has a unique best approximation in \(M\).

%%%%%%%%%%%%%%%%%%%%%%%%%%%%%%%%%%%%%%%%%%%%%%%%%%%%%%%%%%%%
\subsection{Orthogonal Projection}
\label{subsec:approximation-orthogonal-projection}
%%%%%%%%%%%%%%%%%%%%%%%%%%%%%%%%%%%%%%%%%%%%%%%%%%%%%%%%%%%%

Let \(H\) be a Hilbert space and let \(M\subseteq H\) be a closed
subspace.

The best approximation \(P_Mf\) is characterized by

\[
\boxed{
f-P_Mf
\perp
M.
}
\]

Thus

\[
\boxed{
\langle
f-P_Mf,
g
\rangle
=
0
}
\]

for every

\[
g\in M.
\]

This is the projection theorem.

%%%%%%%%%%%%%%%%%%%%%%%%%%%%%%%%%%%%%%%%%%%%%%%%%%%%%%%%%%%%
\subsection{Finite-Dimensional Least Squares}
\label{subsec:finite-dimensional-least-squares}
%%%%%%%%%%%%%%%%%%%%%%%%%%%%%%%%%%%%%%%%%%%%%%%%%%%%%%%%%%%%

Suppose

\[
M
=
\operatorname{span}
\{
\phi_1,\ldots,\phi_m
\}.
\]

Seek

\[
p
=
\sum_{j=1}^{m}
c_j\phi_j.
\]

The orthogonality conditions are

\[
\left\langle
f-
\sum_{j=1}^{m}c_j\phi_j,
\phi_k
\right\rangle
=
0.
\]

Therefore,

\[
\boxed{
\sum_{j=1}^{m}
c_j
\langle\phi_j,\phi_k\rangle
=
\langle f,\phi_k\rangle.
}
\]

These are the normal equations.

%%%%%%%%%%%%%%%%%%%%%%%%%%%%%%%%%%%%%%%%%%%%%%%%%%%%%%%%%%%%
\subsection{Orthogonal Bases}
\label{subsec:approximation-orthogonal-bases}
%%%%%%%%%%%%%%%%%%%%%%%%%%%%%%%%%%%%%%%%%%%%%%%%%%%%%%%%%%%%

If

\[
\phi_1,\ldots,\phi_m
\]

are orthogonal, then the coefficients simplify to

\[
\boxed{
c_j
=
\frac{\langle f,\phi_j\rangle}
{\langle\phi_j,\phi_j\rangle}.
}
\]

If the basis is orthonormal,

\[
\boxed{
c_j
=
\langle f,\phi_j\rangle.
}
\]

This is the same projection principle underlying Fourier series.

%%%%%%%%%%%%%%%%%%%%%%%%%%%%%%%%%%%%%%%%%%%%%%%%%%%%%%%%%%%%
\subsection{Polynomial Approximation}
\label{subsec:polynomial-approximation}
%%%%%%%%%%%%%%%%%%%%%%%%%%%%%%%%%%%%%%%%%%%%%%%%%%%%%%%%%%%%

One of the most important approximation families is

\[
\boxed{
\mathcal{P}_n
=
\{
p:
p\text{ is a polynomial of degree at most }n
\}.
}
\]

The basic question is:

\[
\boxed{
\text{How well can a function be approximated by polynomials?}
}
\]

The answer begins with the Weierstrass approximation theorem.

%%%%%%%%%%%%%%%%%%%%%%%%%%%%%%%%%%%%%%%%%%%%%%%%%%%%%%%%%%%%
\subsection{Weierstrass Approximation Theorem}
\label{subsec:weierstrass-approximation}
%%%%%%%%%%%%%%%%%%%%%%%%%%%%%%%%%%%%%%%%%%%%%%%%%%%%%%%%%%%%

\begin{theorem}[Weierstrass Approximation Theorem]
Let

\[
f\in C([a,b]).
\]

For every

\[
\varepsilon>0,
\]

there exists a polynomial \(p\) such that

\[
\boxed{
\|f-p\|_\infty<\varepsilon.
}
\]
\end{theorem}

Equivalently, the polynomials are dense in

\[
\boxed{
C([a,b])
}
\]

with respect to the uniform norm.

%%%%%%%%%%%%%%%%%%%%%%%%%%%%%%%%%%%%%%%%%%%%%%%%%%%%%%%%%%%%
\subsection{Meaning of Density}
\label{subsec:approximation-density}
%%%%%%%%%%%%%%%%%%%%%%%%%%%%%%%%%%%%%%%%%%%%%%%%%%%%%%%%%%%%

A subset \(A\) of a normed space \(X\) is dense in \(X\) if every element
of \(X\) can be approximated arbitrarily closely by elements of \(A\).

Symbolically,

\[
\boxed{
\overline{A}=X.
}
\]

Thus the Weierstrass theorem states

\[
\boxed{
\overline{
\bigcup_{n=0}^{\infty}\mathcal{P}_n
}^{\,\|\cdot\|_\infty}
=
C([a,b]).
}
\]

This is an existence theorem.

It does not by itself identify the best approximating polynomial or the
rate of convergence.

%%%%%%%%%%%%%%%%%%%%%%%%%%%%%%%%%%%%%%%%%%%%%%%%%%%%%%%%%%%%
\subsection{Bernstein Polynomials}
\label{subsec:bernstein-polynomials}
%%%%%%%%%%%%%%%%%%%%%%%%%%%%%%%%%%%%%%%%%%%%%%%%%%%%%%%%%%%%

Bernstein polynomials provide a constructive proof of the Weierstrass
theorem on \([0,1]\).

For

\[
f\in C([0,1]),
\]

define

\[
\boxed{
B_nf(x)
=
\sum_{k=0}^{n}
f\left(\frac{k}{n}\right)
\binom{n}{k}
x^k(1-x)^{n-k}.
}
\]

Then \(B_nf\) is a polynomial of degree at most \(n\).

%%%%%%%%%%%%%%%%%%%%%%%%%%%%%%%%%%%%%%%%%%%%%%%%%%%%%%%%%%%%
\subsection{Bernstein Convergence}
\label{subsec:bernstein-convergence}
%%%%%%%%%%%%%%%%%%%%%%%%%%%%%%%%%%%%%%%%%%%%%%%%%%%%%%%%%%%%

\begin{theorem}[Bernstein Approximation]
For every

\[
f\in C([0,1]),
\]

the Bernstein polynomials satisfy

\[
\boxed{
\|B_nf-f\|_\infty\to0.
}
\]
\end{theorem}

Thus Bernstein polynomials give an explicit sequence of polynomial
approximations.

%%%%%%%%%%%%%%%%%%%%%%%%%%%%%%%%%%%%%%%%%%%%%%%%%%%%%%%%%%%%
\subsection{Probabilistic Interpretation of Bernstein Polynomials}
\label{subsec:bernstein-probabilistic}
%%%%%%%%%%%%%%%%%%%%%%%%%%%%%%%%%%%%%%%%%%%%%%%%%%%%%%%%%%%%

Let

\[
X_n
\sim
\operatorname{Binomial}(n,x).
\]

Then

\[
\boxed{
B_nf(x)
=
\mathbb{E}
\left[
f\left(\frac{X_n}{n}\right)
\right].
}
\]

Since

\[
\frac{X_n}{n}
\]

concentrates near \(x\) for large \(n\), continuity suggests

\[
f\left(\frac{X_n}{n}\right)
\approx
f(x).
\]

This gives an intuitive probabilistic explanation for Bernstein
approximation.

%%%%%%%%%%%%%%%%%%%%%%%%%%%%%%%%%%%%%%%%%%%%%%%%%%%%%%%%%%%%
\subsection{Worked Example: Bernstein Approximation}
\label{subsec:worked-bernstein}
%%%%%%%%%%%%%%%%%%%%%%%%%%%%%%%%%%%%%%%%%%%%%%%%%%%%%%%%%%%%

\begin{workedexamplebox}{Bernstein Polynomial for \(f(x)=x^2\)}

For \(n=2\),

\[
B_2f(x)
=
\sum_{k=0}^{2}
f\left(\frac{k}{2}\right)
\binom{2}{k}
x^k(1-x)^{2-k}.
\]

Since

\[
f(0)=0,
\qquad
f\left(\frac12\right)=\frac14,
\qquad
f(1)=1,
\]

we obtain

\[
B_2f(x)
=
\frac14
\binom21
x(1-x)
+
x^2.
\]

Therefore,

\[
B_2f(x)
=
\frac12x(1-x)+x^2.
\]

Hence

\begin{answerbox}
\[
\boxed{
B_2f(x)
=
\frac12x+\frac12x^2.
}
\]
\end{answerbox}

\end{workedexamplebox}

%%%%%%%%%%%%%%%%%%%%%%%%%%%%%%%%%%%%%%%%%%%%%%%%%%%%%%%%%%%%
\subsection{Positive Linear Operators}
\label{subsec:positive-linear-operators}
%%%%%%%%%%%%%%%%%%%%%%%%%%%%%%%%%%%%%%%%%%%%%%%%%%%%%%%%%%%%

Bernstein approximation is an example of approximation by positive
linear operators.

An operator \(L\) is positive if

\[
f\geq0
\]

implies

\[
\boxed{
Lf\geq0.
}
\]

The Bernstein operator satisfies

\[
\boxed{
B_n(af+bg)
=
aB_nf+bB_ng
}
\]

and preserves nonnegativity.

Positive operators play an important role in constructive approximation
theory.

%%%%%%%%%%%%%%%%%%%%%%%%%%%%%%%%%%%%%%%%%%%%%%%%%%%%%%%%%%%%
\subsection{Korovkin-Type Principle}
\label{subsec:korovkin-principle}
%%%%%%%%%%%%%%%%%%%%%%%%%%%%%%%%%%%%%%%%%%%%%%%%%%%%%%%%%%%%

A remarkable feature of positive approximation operators is that their
convergence can sometimes be checked on only a few simple functions.

For positive linear operators on \(C([a,b])\), a classical
Korovkin-type theorem reduces uniform convergence to checking suitable
test functions corresponding to

\[
\boxed{
1,
\qquad
x,
\qquad
x^2.
}
\]

This greatly simplifies convergence proofs for many approximation
schemes.

%%%%%%%%%%%%%%%%%%%%%%%%%%%%%%%%%%%%%%%%%%%%%%%%%%%%%%%%%%%%
\subsection{Stone--Weierstrass Theorem}
\label{subsec:stone-weierstrass}
%%%%%%%%%%%%%%%%%%%%%%%%%%%%%%%%%%%%%%%%%%%%%%%%%%%%%%%%%%%%

The Weierstrass theorem extends far beyond intervals.

A standard real form is the following.

\begin{theorem}[Stone--Weierstrass Theorem]
Let \(K\) be a compact Hausdorff space and let \(A\) be a subalgebra of
\(C(K,\R)\).

If \(A\) contains the constant functions and separates points of \(K\),
then \(A\) is dense in \(C(K,\R)\) in the uniform norm.
\end{theorem}

To separate points means that for any distinct \(x,y\in K\), there exists

\[
f\in A
\]

such that

\[
\boxed{
f(x)\neq f(y).
}
\]

%%%%%%%%%%%%%%%%%%%%%%%%%%%%%%%%%%%%%%%%%%%%%%%%%%%%%%%%%%%%
\subsection{Why Stone--Weierstrass Matters}
\label{subsec:why-stone-weierstrass}
%%%%%%%%%%%%%%%%%%%%%%%%%%%%%%%%%%%%%%%%%%%%%%%%%%%%%%%%%%%%

Stone--Weierstrass gives a general recipe for proving density.

Rather than explicitly constructing approximations, one checks
algebraic properties.

It explains why many familiar approximation systems are dense, including
appropriate families of

\[
\boxed{
\text{polynomials},
\quad
\text{trigonometric polynomials},
\quad
\text{algebra-generated function families}.
}
\]

%%%%%%%%%%%%%%%%%%%%%%%%%%%%%%%%%%%%%%%%%%%%%%%%%%%%%%%%%%%%
\subsection{Trigonometric Approximation}
\label{subsec:trigonometric-approximation}
%%%%%%%%%%%%%%%%%%%%%%%%%%%%%%%%%%%%%%%%%%%%%%%%%%%%%%%%%%%%

A trigonometric polynomial has the form

\[
\boxed{
T_n(x)
=
a_0
+
\sum_{k=1}^{n}
\left(
a_k\cos(kx)
+
b_k\sin(kx)
\right).
}
\]

Equivalently,

\[
\boxed{
T_n(x)
=
\sum_{k=-n}^{n}
c_ke^{ikx}.
}
\]

Trigonometric polynomials form the natural approximation family for
periodic functions.

%%%%%%%%%%%%%%%%%%%%%%%%%%%%%%%%%%%%%%%%%%%%%%%%%%%%%%%%%%%%
\subsection{Fourier Approximation}
\label{subsec:fourier-as-approximation}
%%%%%%%%%%%%%%%%%%%%%%%%%%%%%%%%%%%%%%%%%%%%%%%%%%%%%%%%%%%%

In \(L^2\), the Fourier partial sum is an orthogonal projection onto a
finite-dimensional trigonometric subspace.

Therefore,

\[
\boxed{
S_nf
}
\]

is the best \(L^2\) approximation to \(f\) from that subspace.

This gives a geometric interpretation of Fourier series:

\[
\boxed{
\text{Fourier approximation}
=
\text{Hilbert-space projection}.
}
\]

%%%%%%%%%%%%%%%%%%%%%%%%%%%%%%%%%%%%%%%%%%%%%%%%%%%%%%%%%%%%
\subsection{Orthogonal Polynomials}
\label{subsec:orthogonal-polynomials}
%%%%%%%%%%%%%%%%%%%%%%%%%%%%%%%%%%%%%%%%%%%%%%%%%%%%%%%%%%%%

Let \(w(x)\geq0\) be a weight function on an interval \(I\).

Define

\[
\boxed{
\langle f,g\rangle_w
=
\int_I
f(x)g(x)w(x)\,dx.
}
\]

A sequence of polynomials

\[
p_0,p_1,p_2,\ldots
\]

is orthogonal if

\[
\boxed{
\langle p_j,p_k\rangle_w=0
\qquad
(j\neq k).
}
\]

Orthogonal polynomial systems provide efficient bases for polynomial
approximation.

%%%%%%%%%%%%%%%%%%%%%%%%%%%%%%%%%%%%%%%%%%%%%%%%%%%%%%%%%%%%
\subsection{Important Orthogonal Polynomial Families}
\label{subsec:orthogonal-polynomial-families}
%%%%%%%%%%%%%%%%%%%%%%%%%%%%%%%%%%%%%%%%%%%%%%%%%%%%%%%%%%%%

Classical examples include

\[
\boxed{
\text{Legendre polynomials}
}
\]

on \([-1,1]\) with weight \(1\),

\[
\boxed{
\text{Chebyshev polynomials}
}
\]

with weight

\[
\frac{1}{\sqrt{1-x^2}},
\]

\[
\boxed{
\text{Hermite polynomials}
}
\]

on \(\R\) with Gaussian-type weight, and

\[
\boxed{
\text{Laguerre polynomials}
}
\]

on \([0,\infty)\) with exponential-type weight.

These systems appear throughout approximation theory, numerical analysis,
differential equations, and mathematical physics.

%%%%%%%%%%%%%%%%%%%%%%%%%%%%%%%%%%%%%%%%%%%%%%%%%%%%%%%%%%%%
\subsection{Legendre Approximation}
\label{subsec:legendre-approximation}
%%%%%%%%%%%%%%%%%%%%%%%%%%%%%%%%%%%%%%%%%%%%%%%%%%%%%%%%%%%%

The Legendre polynomials \(P_n\) satisfy

\[
\boxed{
\int_{-1}^{1}
P_m(x)P_n(x)\,dx
=
0
\qquad
(m\neq n).
}
\]

A function may be approximated by

\[
\boxed{
f(x)
\approx
\sum_{k=0}^{N}
c_kP_k(x).
}
\]

The coefficients are obtained by orthogonal projection:

\[
\boxed{
c_k
=
\frac{
\int_{-1}^{1}
f(x)P_k(x)\,dx
}{
\int_{-1}^{1}
P_k(x)^2\,dx
}.
}
\]

%%%%%%%%%%%%%%%%%%%%%%%%%%%%%%%%%%%%%%%%%%%%%%%%%%%%%%%%%%%%
\subsection{Chebyshev Polynomials}
\label{subsec:chebyshev-polynomials}
%%%%%%%%%%%%%%%%%%%%%%%%%%%%%%%%%%%%%%%%%%%%%%%%%%%%%%%%%%%%

The Chebyshev polynomial of the first kind is defined by

\[
\boxed{
T_n(x)
=
\cos(n\arccos x),
\qquad
-1\leq x\leq1.
}
\]

Equivalently,

\[
\boxed{
T_n(\cos\theta)
=
\cos(n\theta).
}
\]

The first few are

\[
\boxed{
T_0(x)=1,
}
\]

\[
\boxed{
T_1(x)=x,
}
\]

\[
\boxed{
T_2(x)=2x^2-1,
}
\]

and

\[
\boxed{
T_3(x)=4x^3-3x.
}
\]

%%%%%%%%%%%%%%%%%%%%%%%%%%%%%%%%%%%%%%%%%%%%%%%%%%%%%%%%%%%%
\subsection{Chebyshev Recurrence}
\label{subsec:chebyshev-recurrence}
%%%%%%%%%%%%%%%%%%%%%%%%%%%%%%%%%%%%%%%%%%%%%%%%%%%%%%%%%%%%

Chebyshev polynomials satisfy

\[
\boxed{
T_{n+1}(x)
=
2xT_n(x)-T_{n-1}(x).
}
\]

They also satisfy

\[
\boxed{
|T_n(x)|\leq1
}
\]

for

\[
x\in[-1,1].
\]

Their oscillatory structure makes them central to uniform approximation.

%%%%%%%%%%%%%%%%%%%%%%%%%%%%%%%%%%%%%%%%%%%%%%%%%%%%%%%%%%%%
\subsection{Chebyshev Nodes}
\label{subsec:chebyshev-nodes}
%%%%%%%%%%%%%%%%%%%%%%%%%%%%%%%%%%%%%%%%%%%%%%%%%%%%%%%%%%%%

The zeros of \(T_n\) are

\[
\boxed{
x_k
=
\cos
\left(
\frac{(2k-1)\pi}{2n}
\right),
\qquad
k=1,\ldots,n.
}
\]

These are called Chebyshev nodes.

They cluster near the endpoints of \([-1,1]\), which helps control large
interpolation errors near the boundaries.

%%%%%%%%%%%%%%%%%%%%%%%%%%%%%%%%%%%%%%%%%%%%%%%%%%%%%%%%%%%%
\subsection{Runge Phenomenon}
\label{subsec:runge-phenomenon}
%%%%%%%%%%%%%%%%%%%%%%%%%%%%%%%%%%%%%%%%%%%%%%%%%%%%%%%%%%%%

High-degree polynomial interpolation at equally spaced points can behave
poorly near interval endpoints.

This is called the Runge phenomenon.

Increasing the number of equally spaced interpolation points can
therefore make the approximation worse.

Chebyshev nodes dramatically reduce this instability.

This illustrates an important distinction:

\[
\boxed{
\text{high polynomial degree}
\not\Rightarrow
\text{good approximation automatically}.
}
\]

%%%%%%%%%%%%%%%%%%%%%%%%%%%%%%%%%%%%%%%%%%%%%%%%%%%%%%%%%%%%
\subsection{Uniform Best Approximation}
\label{subsec:uniform-best-approximation}
%%%%%%%%%%%%%%%%%%%%%%%%%%%%%%%%%%%%%%%%%%%%%%%%%%%%%%%%%%%%

For

\[
f\in C([a,b]),
\]

define the degree-\(n\) best uniform approximation error by

\[
\boxed{
E_n(f)_\infty
=
\inf_{p\in\mathcal{P}_n}
\|f-p\|_\infty.
}
\]

A polynomial attaining this infimum is called a minimax or best uniform
approximating polynomial.

%%%%%%%%%%%%%%%%%%%%%%%%%%%%%%%%%%%%%%%%%%%%%%%%%%%%%%%%%%%%
\subsection{Minimax Approximation}
\label{subsec:minimax-approximation}
%%%%%%%%%%%%%%%%%%%%%%%%%%%%%%%%%%%%%%%%%%%%%%%%%%%%%%%%%%%%

Minimax approximation minimizes the largest absolute error:

\[
\boxed{
\min_{p\in\mathcal{P}_n}
\max_{x\in[a,b]}
|f(x)-p(x)|.
}
\]

This differs from least-squares approximation, which minimizes an
integrated squared error.

Thus

\[
\boxed{
\text{least squares}
\longrightarrow
L^2,
}
\]

while

\[
\boxed{
\text{minimax}
\longrightarrow
L^\infty.
}
\]

%%%%%%%%%%%%%%%%%%%%%%%%%%%%%%%%%%%%%%%%%%%%%%%%%%%%%%%%%%%%
\subsection{Equioscillation Principle}
\label{subsec:equioscillation-principle}
%%%%%%%%%%%%%%%%%%%%%%%%%%%%%%%%%%%%%%%%%%%%%%%%%%%%%%%%%%%%

A central characterization of best uniform polynomial approximation is
the alternation theorem.

\begin{theorem}[Chebyshev Equioscillation Theorem]
Let \(f\in C([a,b])\), and let \(p^\ast\in\mathcal{P}_n\).

Then \(p^\ast\) is the unique best uniform approximation to \(f\) from
\(\mathcal{P}_n\) if and only if there exist at least \(n+2\) points

\[
a\leq x_0<x_1<\cdots<x_{n+1}\leq b
\]

such that

\[
\boxed{
f(x_j)-p^\ast(x_j)
=
(-1)^j\sigma E
}
\]

for some \(\sigma\in\{-1,1\}\), where

\[
E=\|f-p^\ast\|_\infty.
\]
\end{theorem}

Thus the optimal error alternates between equal positive and negative
extremes.

%%%%%%%%%%%%%%%%%%%%%%%%%%%%%%%%%%%%%%%%%%%%%%%%%%%%%%%%%%%%
\subsection{Worked Example: Best Constant Approximation}
\label{subsec:worked-best-constant}
%%%%%%%%%%%%%%%%%%%%%%%%%%%%%%%%%%%%%%%%%%%%%%%%%%%%%%%%%%%%

\begin{workedexamplebox}{Best Constant Approximation to \(f(x)=x^2\) on \([-1,1]\)}

Seek a constant \(c\) minimizing

\[
\|x^2-c\|_\infty.
\]

Since

\[
0\leq x^2\leq1,
\]

the maximum error is minimized by placing \(c\) halfway between the
minimum and maximum values.

Thus

\[
c=\frac12.
\]

The errors at \(x=0\) and \(x=\pm1\) are

\[
-\frac12
\]

and

\[
\frac12,
\]

respectively.

Therefore,

\begin{answerbox}
\[
\boxed{
p^\ast(x)=\frac12,
\qquad
E_0(f)_\infty=\frac12.
}
\]
\end{answerbox}

\end{workedexamplebox}

%%%%%%%%%%%%%%%%%%%%%%%%%%%%%%%%%%%%%%%%%%%%%%%%%%%%%%%%%%%%
\subsection{Approximation Rate}
\label{subsec:approximation-rate}
%%%%%%%%%%%%%%%%%%%%%%%%%%%%%%%%%%%%%%%%%%%%%%%%%%%%%%%%%%%%

Density answers whether approximation is possible.

A deeper question asks

\[
\boxed{
\text{How fast does the error decrease?}
}
\]

For polynomial approximation, one studies quantities such as

\[
\boxed{
E_n(f)_\infty
=
\inf_{p\in\mathcal{P}_n}
\|f-p\|_\infty.
}
\]

The behavior of \(E_n(f)_\infty\) as

\[
n\to\infty
\]

is closely related to the smoothness of \(f\).

%%%%%%%%%%%%%%%%%%%%%%%%%%%%%%%%%%%%%%%%%%%%%%%%%%%%%%%%%%%%
\subsection{Smoothness and Approximation}
\label{subsec:smoothness-approximation}
%%%%%%%%%%%%%%%%%%%%%%%%%%%%%%%%%%%%%%%%%%%%%%%%%%%%%%%%%%%%

As a general principle,

\[
\boxed{
\text{greater smoothness}
\Longrightarrow
\text{faster approximation}.
}
\]

Functions with only limited regularity generally have slower polynomial
or Fourier approximation rates.

Functions possessing many derivatives typically have faster algebraic
decay of approximation error.

Analytic functions may exhibit geometric or exponential-type convergence
under suitable approximation schemes.

%%%%%%%%%%%%%%%%%%%%%%%%%%%%%%%%%%%%%%%%%%%%%%%%%%%%%%%%%%%%
\subsection{Modulus of Continuity}
\label{subsec:modulus-of-continuity}
%%%%%%%%%%%%%%%%%%%%%%%%%%%%%%%%%%%%%%%%%%%%%%%%%%%%%%%%%%%%

\begin{definitionbox}{Modulus of Continuity}
For a continuous function \(f\), define

\[
\boxed{
\omega(f,\delta)
=
\sup_{\substack{x,y\in[a,b]\\|x-y|\leq\delta}}
|f(x)-f(y)|.
}
\]
\end{definitionbox}

The modulus of continuity measures how much \(f\) can change over
distances no greater than \(\delta\).

For a uniformly continuous function,

\[
\boxed{
\omega(f,\delta)\to0
\qquad
\text{as }\delta\to0^+.
}
\]

%%%%%%%%%%%%%%%%%%%%%%%%%%%%%%%%%%%%%%%%%%%%%%%%%%%%%%%%%%%%
\subsection{Jackson-Type Estimates}
\label{subsec:jackson-estimates}
%%%%%%%%%%%%%%%%%%%%%%%%%%%%%%%%%%%%%%%%%%%%%%%%%%%%%%%%%%%%

Jackson-type theorems connect smoothness directly with best approximation
error.

A representative form is

\[
\boxed{
E_n(f)_\infty
\leq
C\,
\omega
\left(
f,\frac{1}{n}
\right),
}
\]

where \(C\) is independent of \(f\) and \(n\).

More refined versions involve higher derivatives and higher-order moduli
of smoothness.

These are called direct approximation theorems.

%%%%%%%%%%%%%%%%%%%%%%%%%%%%%%%%%%%%%%%%%%%%%%%%%%%%%%%%%%%%
\subsection{Inverse Approximation Theorems}
\label{subsec:inverse-approximation-theorems}
%%%%%%%%%%%%%%%%%%%%%%%%%%%%%%%%%%%%%%%%%%%%%%%%%%%%%%%%%%%%

Direct theorems say

\[
\boxed{
\text{smoothness}
\Longrightarrow
\text{fast approximation}.
}
\]

Inverse theorems ask whether

\[
\boxed{
\text{fast approximation}
\Longrightarrow
\text{smoothness}.
}
\]

Under suitable hypotheses, the answer is yes.

Together, direct and inverse theorems establish deep equivalences between
regularity and approximability.

%%%%%%%%%%%%%%%%%%%%%%%%%%%%%%%%%%%%%%%%%%%%%%%%%%%%%%%%%%%%
\subsection{Piecewise Polynomial Approximation}
\label{subsec:piecewise-polynomial-approximation}
%%%%%%%%%%%%%%%%%%%%%%%%%%%%%%%%%%%%%%%%%%%%%%%%%%%%%%%%%%%%

Instead of using one high-degree polynomial on an entire interval, one
can divide the interval into smaller pieces and use low-degree
polynomials locally.

Let

\[
a=x_0<x_1<\cdots<x_N=b.
\]

On each subinterval

\[
[x_{j-1},x_j],
\]

one approximates \(f\) by a polynomial.

This leads to piecewise polynomial approximation.

%%%%%%%%%%%%%%%%%%%%%%%%%%%%%%%%%%%%%%%%%%%%%%%%%%%%%%%%%%%%
\subsection{Splines}
\label{subsec:splines}
%%%%%%%%%%%%%%%%%%%%%%%%%%%%%%%%%%%%%%%%%%%%%%%%%%%%%%%%%%%%

\begin{definitionbox}{Spline}
A spline is a piecewise polynomial function whose polynomial pieces are
joined with prescribed smoothness conditions at specified points called
knots.
\end{definitionbox}

If the knots are

\[
x_0<x_1<\cdots<x_N,
\]

then a spline is polynomial on each interval

\[
[x_{j-1},x_j].
\]

The degree and continuity conditions determine the spline class.

%%%%%%%%%%%%%%%%%%%%%%%%%%%%%%%%%%%%%%%%%%%%%%%%%%%%%%%%%%%%
\subsection{Cubic Splines}
\label{subsec:cubic-splines}
%%%%%%%%%%%%%%%%%%%%%%%%%%%%%%%%%%%%%%%%%%%%%%%%%%%%%%%%%%%%

A cubic spline is piecewise cubic and typically satisfies

\[
\boxed{
S\in C^2([a,b]).
}
\]

Thus

\[
S,
\qquad
S',
\qquad
S''
\]

are continuous across the interior knots.

Cubic splines are widely used because they provide a strong balance
between

\[
\boxed{
\text{smoothness},
\quad
\text{accuracy},
\quad
\text{computational simplicity}.
}
\]

%%%%%%%%%%%%%%%%%%%%%%%%%%%%%%%%%%%%%%%%%%%%%%%%%%%%%%%%%%%%
\subsection{Why Splines Work Well}
\label{subsec:why-splines-work}
%%%%%%%%%%%%%%%%%%%%%%%%%%%%%%%%%%%%%%%%%%%%%%%%%%%%%%%%%%%%

A single high-degree polynomial has global behavior: changing one
coefficient affects the entire interval.

Splines are local.

Changing data in one region primarily affects nearby pieces.

This gives splines improved stability and flexibility.

It also avoids many difficulties associated with very high-degree global
polynomial interpolation.

%%%%%%%%%%%%%%%%%%%%%%%%%%%%%%%%%%%%%%%%%%%%%%%%%%%%%%%%%%%%
\subsection{B-Splines}
\label{subsec:b-splines}
%%%%%%%%%%%%%%%%%%%%%%%%%%%%%%%%%%%%%%%%%%%%%%%%%%%%%%%%%%%%

B-splines provide a particularly useful basis for spline spaces.

Their important properties include

\[
\boxed{
\text{local support},
}
\]

\[
\boxed{
\text{nonnegativity},
}
\]

and

\[
\boxed{
\text{stable numerical evaluation}.
}
\]

Because each basis function is supported only on a small number of knot
intervals, B-spline representations are computationally efficient.

%%%%%%%%%%%%%%%%%%%%%%%%%%%%%%%%%%%%%%%%%%%%%%%%%%%%%%%%%%%%
\subsection{Rational Approximation}
\label{subsec:rational-approximation}
%%%%%%%%%%%%%%%%%%%%%%%%%%%%%%%%%%%%%%%%%%%%%%%%%%%%%%%%%%%%

Instead of approximating by a polynomial, one may use a rational function

\[
\boxed{
R(x)
=
\frac{p(x)}{q(x)}.
}
\]

Rational functions can approximate functions with poles, sharp
transitions, or other difficult behavior much more efficiently than
polynomials in many settings.

A rational approximation problem may seek

\[
\boxed{
\min_{R}
\|f-R\|.
}
\]

%%%%%%%%%%%%%%%%%%%%%%%%%%%%%%%%%%%%%%%%%%%%%%%%%%%%%%%%%%%%
\subsection{Pad\'e Approximation}
\label{subsec:pade-approximation}
%%%%%%%%%%%%%%%%%%%%%%%%%%%%%%%%%%%%%%%%%%%%%%%%%%%%%%%%%%%%

A Pad\'e approximant is a rational function

\[
\boxed{
R_{m,n}(x)
=
\frac{P_m(x)}{Q_n(x)}
}
\]

chosen so that its power-series expansion agrees with that of \(f\) to
as high an order as possible near a specified point.

Unlike Taylor polynomials, Pad\'e approximants can represent pole-like
behavior.

This often gives much better approximation outside a small neighborhood
of the expansion point.

%%%%%%%%%%%%%%%%%%%%%%%%%%%%%%%%%%%%%%%%%%%%%%%%%%%%%%%%%%%%
\subsection{Taylor Approximation Versus Best Approximation}
\label{subsec:taylor-versus-best-approximation}
%%%%%%%%%%%%%%%%%%%%%%%%%%%%%%%%%%%%%%%%%%%%%%%%%%%%%%%%%%%%

A Taylor polynomial is determined by matching derivatives at one point:

\[
\boxed{
p^{(k)}(a)=f^{(k)}(a).
}
\]

A best approximation is determined by minimizing an error over an entire
domain.

Therefore a Taylor polynomial need not be the best approximation in
either

\[
L^2
\]

or

\[
L^\infty.
\]

The two methods solve different problems.

%%%%%%%%%%%%%%%%%%%%%%%%%%%%%%%%%%%%%%%%%%%%%%%%%%%%%%%%%%%%
\subsection{Approximation by Convolution}
\label{subsec:approximation-convolution}
%%%%%%%%%%%%%%%%%%%%%%%%%%%%%%%%%%%%%%%%%%%%%%%%%%%%%%%%%%%%

Approximate identities provide another major approximation mechanism.

Let

\[
\phi_\varepsilon(x)
=
\varepsilon^{-n}
\phi(x/\varepsilon),
\]

where

\[
\int_{\R^n}\phi(x)\,dx=1.
\]

Then

\[
\boxed{
f_\varepsilon
=
f*\phi_\varepsilon
}
\]

often satisfies

\[
\boxed{
f_\varepsilon\to f
}
\]

in an appropriate norm or pointwise sense.

If \(\phi\) is smooth, then \(f_\varepsilon\) is smoother than \(f\).

%%%%%%%%%%%%%%%%%%%%%%%%%%%%%%%%%%%%%%%%%%%%%%%%%%%%%%%%%%%%
\subsection{Mollification}
\label{subsec:approximation-mollification}
%%%%%%%%%%%%%%%%%%%%%%%%%%%%%%%%%%%%%%%%%%%%%%%%%%%%%%%%%%%%

When

\[
\phi\in C_c^\infty(\R^n),
\]

the functions

\[
f_\varepsilon=f*\phi_\varepsilon
\]

are called mollifications of \(f\).

Mollification is used to approximate rough functions by smooth ones.

For example, under standard conditions,

\[
\boxed{
f_\varepsilon\to f
\quad
\text{in }L^p.
}
\]

This technique is fundamental in real analysis, distribution theory, and
PDEs.

%%%%%%%%%%%%%%%%%%%%%%%%%%%%%%%%%%%%%%%%%%%%%%%%%%%%%%%%%%%%
\subsection{Approximation in Sobolev Spaces}
\label{subsec:approximation-sobolev}
%%%%%%%%%%%%%%%%%%%%%%%%%%%%%%%%%%%%%%%%%%%%%%%%%%%%%%%%%%%%

Approximation may involve both functions and their derivatives.

For example, convergence in

\[
W^{1,p}(\Omega)
\]

requires

\[
\boxed{
\|f_n-f\|_{L^p}
+
\|\nabla f_n-\nabla f\|_{L^p}
\to0.
}
\]

Under appropriate conditions, smooth functions are dense in Sobolev
spaces.

This density is essential for weak formulations of differential
equations.

%%%%%%%%%%%%%%%%%%%%%%%%%%%%%%%%%%%%%%%%%%%%%%%%%%%%%%%%%%%%
\subsection{Approximation in Hilbert Spaces}
\label{subsec:approximation-hilbert-spaces}
%%%%%%%%%%%%%%%%%%%%%%%%%%%%%%%%%%%%%%%%%%%%%%%%%%%%%%%%%%%%

Let

\[
V_1\subseteq V_2\subseteq\cdots
\]

be finite-dimensional subspaces of a Hilbert space \(H\).

If

\[
\boxed{
\overline{
\bigcup_{n=1}^{\infty}V_n
}
=
H,
}
\]

then orthogonal projections

\[
P_{V_n}f
\]

approximate every \(f\in H\):

\[
\boxed{
\|f-P_{V_n}f\|\to0.
}
\]

This abstract principle underlies many numerical approximation methods.

%%%%%%%%%%%%%%%%%%%%%%%%%%%%%%%%%%%%%%%%%%%%%%%%%%%%%%%%%%%%
\subsection{Galerkin Approximation}
\label{subsec:galerkin-approximation}
%%%%%%%%%%%%%%%%%%%%%%%%%%%%%%%%%%%%%%%%%%%%%%%%%%%%%%%%%%%%

Suppose an equation is written abstractly as

\[
Au=f.
\]

Choose a finite-dimensional space

\[
V_n
=
\operatorname{span}
\{
\phi_1,\ldots,\phi_n
\}.
\]

Seek

\[
u_n\in V_n
\]

such that the residual is orthogonal to the approximation space.

Schematically,

\[
\boxed{
\langle Au_n-f,v\rangle=0
\qquad
\text{for all }v\in V_n.
}
\]

This is the Galerkin principle.

%%%%%%%%%%%%%%%%%%%%%%%%%%%%%%%%%%%%%%%%%%%%%%%%%%%%%%%%%%%%
\subsection{Finite Element Connection}
\label{subsec:finite-element-connection}
%%%%%%%%%%%%%%%%%%%%%%%%%%%%%%%%%%%%%%%%%%%%%%%%%%%%%%%%%%%%

Finite element methods approximate solutions of differential equations
using finite-dimensional spaces of piecewise polynomial functions.

Thus they combine

\[
\boxed{
\text{approximation theory},
\quad
\text{variational methods},
\quad
\text{linear algebra},
\quad
\text{numerical analysis}.
}
\]

Approximation theory provides the estimates needed to determine how the
numerical solution approaches the true solution.

%%%%%%%%%%%%%%%%%%%%%%%%%%%%%%%%%%%%%%%%%%%%%%%%%%%%%%%%%%%%
\subsection{Approximation and Compactness}
\label{subsec:approximation-compactness}
%%%%%%%%%%%%%%%%%%%%%%%%%%%%%%%%%%%%%%%%%%%%%%%%%%%%%%%%%%%%

Finite-dimensional approximation is closely connected with compactness.

A compact operator can often be approximated in operator norm by
finite-rank operators under suitable hypotheses.

Approximation numbers and related quantities measure how effectively an
operator can be replaced by low-dimensional operators.

This connects Approximation Theory with Functional Analysis and Operator
Theory.

%%%%%%%%%%%%%%%%%%%%%%%%%%%%%%%%%%%%%%%%%%%%%%%%%%%%%%%%%%%%
\subsection{Nonlinear Approximation}
\label{subsec:nonlinear-approximation}
%%%%%%%%%%%%%%%%%%%%%%%%%%%%%%%%%%%%%%%%%%%%%%%%%%%%%%%%%%%%

Not all approximation families are linear spaces.

In nonlinear approximation, the allowed approximants may depend on
parameters nonlinearly.

Examples include

\[
\boxed{
\text{rational functions},
\quad
\text{adaptive splines},
\quad
\text{sparse expansions},
\quad
\text{neural-network approximation}.
}
\]

The approximation problem becomes an optimization problem over a
nonlinear family.

%%%%%%%%%%%%%%%%%%%%%%%%%%%%%%%%%%%%%%%%%%%%%%%%%%%%%%%%%%%%
\subsection{Sparse Approximation}
\label{subsec:sparse-approximation}
%%%%%%%%%%%%%%%%%%%%%%%%%%%%%%%%%%%%%%%%%%%%%%%%%%%%%%%%%%%%

Suppose a function has an expansion

\[
f
=
\sum_{k}
c_k\phi_k.
\]

Instead of retaining every coefficient, sparse approximation seeks a
small subset of important terms.

An \(m\)-term approximation has the form

\[
\boxed{
f_m
=
\sum_{k\in\Lambda_m}
c_k\phi_k,
}
\]

where

\[
|\Lambda_m|=m.
\]

Wavelet approximation is a major setting in which sparse representations
are effective.

%%%%%%%%%%%%%%%%%%%%%%%%%%%%%%%%%%%%%%%%%%%%%%%%%%%%%%%%%%%%
\subsection{Approximation and Wavelets}
\label{subsec:approximation-wavelets}
%%%%%%%%%%%%%%%%%%%%%%%%%%%%%%%%%%%%%%%%%%%%%%%%%%%%%%%%%%%%

Wavelets provide localized basis functions across multiple scales.

A function may be approximated by retaining only the most significant
wavelet coefficients.

For functions with localized irregularities, this can be much more
efficient than global Fourier approximation.

This is important in

\[
\boxed{
\text{compression},
\quad
\text{denoising},
\quad
\text{adaptive numerical methods}.
}
\]

%%%%%%%%%%%%%%%%%%%%%%%%%%%%%%%%%%%%%%%%%%%%%%%%%%%%%%%%%%%%
\subsection{Universal Approximation Perspective}
\label{subsec:universal-approximation-perspective}
%%%%%%%%%%%%%%%%%%%%%%%%%%%%%%%%%%%%%%%%%%%%%%%%%%%%%%%%%%%%

Many approximation theorems have the same structural form:

\[
\boxed{
\overline{\mathcal{A}}=X.
}
\]

This says that a relatively simple family \(\mathcal{A}\) is rich enough
to approximate every element of a larger function space \(X\).

Polynomial density, trigonometric approximation, spline density, and
certain neural-network approximation results all fit this broad
perspective.

Density, however, says only that approximation is possible.

It does not automatically provide

\[
\boxed{
\text{efficient computation},
\quad
\text{stability},
\quad
\text{fast convergence}.
}
\]

%%%%%%%%%%%%%%%%%%%%%%%%%%%%%%%%%%%%%%%%%%%%%%%%%%%%%%%%%%%%
\subsection{Stability}
\label{subsec:approximation-stability}
%%%%%%%%%%%%%%%%%%%%%%%%%%%%%%%%%%%%%%%%%%%%%%%%%%%%%%%%%%%%

A useful approximation method should not amplify small perturbations
dramatically.

Suppose data \(f\) are replaced by

\[
f+\delta f.
\]

A stable approximation procedure should produce a controlled change in
the approximant.

Thus practical approximation requires balancing

\[
\boxed{
\text{accuracy}
}
\]

and

\[
\boxed{
\text{stability}.
}
\]

This becomes particularly important in numerical computation.

%%%%%%%%%%%%%%%%%%%%%%%%%%%%%%%%%%%%%%%%%%%%%%%%%%%%%%%%%%%%
\subsection{Conditioning and Approximation}
\label{subsec:conditioning-approximation}
%%%%%%%%%%%%%%%%%%%%%%%%%%%%%%%%%%%%%%%%%%%%%%%%%%%%%%%%%%%%

Even if an approximation problem is mathematically well posed, a poor
basis can make its numerical computation unstable.

For example, high-degree monomial bases

\[
1,x,x^2,\ldots,x^n
\]

may lead to badly conditioned systems.

Orthogonal polynomial bases often improve numerical behavior.

Thus the choice of representation is part of the approximation problem.

%%%%%%%%%%%%%%%%%%%%%%%%%%%%%%%%%%%%%%%%%%%%%%%%%%%%%%%%%%%%
\subsection{Common Errors}
\label{subsec:approximation-theory-common-errors}
%%%%%%%%%%%%%%%%%%%%%%%%%%%%%%%%%%%%%%%%%%%%%%%%%%%%%%%%%%%%

\subsubsection{Confusing Approximation with Interpolation}

Interpolation requires exact matching at selected points.

Approximation minimizes or controls an error criterion.

%%%%%%%%%%%%%%%%%%%%%%%%%%%%%%%%%%%%%%%%%%%%%%%%%%%%%%%%%%%%
\subsubsection{Assuming the Taylor Polynomial Is the Best Polynomial}

Taylor polynomials optimize local derivative matching, not global
\(L^2\) or uniform error.

%%%%%%%%%%%%%%%%%%%%%%%%%%%%%%%%%%%%%%%%%%%%%%%%%%%%%%%%%%%%
\subsubsection{Ignoring the Norm}

The phrase ``best approximation'' is incomplete unless the error measure
is specified.

The best \(L^2\) approximation can differ from the best \(L^\infty\)
approximation.

%%%%%%%%%%%%%%%%%%%%%%%%%%%%%%%%%%%%%%%%%%%%%%%%%%%%%%%%%%%%
\subsubsection{Assuming Density Gives a Convergence Rate}

A density theorem proves that arbitrarily accurate approximations exist.

It does not specify how rapidly the error decreases with dimension or
degree.

%%%%%%%%%%%%%%%%%%%%%%%%%%%%%%%%%%%%%%%%%%%%%%%%%%%%%%%%%%%%
\subsubsection{Assuming High Degree Always Improves Numerical Results}

High-degree polynomial interpolation at poorly chosen nodes can become
unstable.

The Runge phenomenon is a standard example.

%%%%%%%%%%%%%%%%%%%%%%%%%%%%%%%%%%%%%%%%%%%%%%%%%%%%%%%%%%%%
\subsubsection{Confusing Orthogonality with Orthonormality}

Orthogonality means

\[
\langle\phi_j,\phi_k\rangle=0
\]

for \(j\neq k\).

Orthonormality additionally requires

\[
\boxed{
\|\phi_j\|=1.
}
\]

%%%%%%%%%%%%%%%%%%%%%%%%%%%%%%%%%%%%%%%%%%%%%%%%%%%%%%%%%%%%
\subsubsection{Assuming Least Squares Minimizes Maximum Error}

Least squares minimizes an \(L^2\)-type quantity.

Minimax approximation minimizes the largest error.

%%%%%%%%%%%%%%%%%%%%%%%%%%%%%%%%%%%%%%%%%%%%%%%%%%%%%%%%%%%%
\subsubsection{Ignoring Smoothness}

Approximation rates depend strongly on regularity.

A nonsmooth function should not generally be expected to converge as
rapidly as an analytic function.

%%%%%%%%%%%%%%%%%%%%%%%%%%%%%%%%%%%%%%%%%%%%%%%%%%%%%%%%%%%%
\subsubsection{Assuming Global Polynomial Approximation Is Always Best}

Piecewise polynomials and splines can be more stable and efficient,
especially for complicated local behavior.

%%%%%%%%%%%%%%%%%%%%%%%%%%%%%%%%%%%%%%%%%%%%%%%%%%%%%%%%%%%%
\subsubsection{Confusing Approximation Error with Numerical Error}

Approximation error comes from restricting the solution to a simpler
family.

Numerical error can additionally include

\[
\boxed{
\text{roundoff},
\quad
\text{discretization},
\quad
\text{iteration error}.
}
\]

These should be distinguished.

%%%%%%%%%%%%%%%%%%%%%%%%%%%%%%%%%%%%%%%%%%%%%%%%%%%%%%%%%%%%
\subsection{Applications}
\label{subsec:approximation-theory-applications}
%%%%%%%%%%%%%%%%%%%%%%%%%%%%%%%%%%%%%%%%%%%%%%%%%%%%%%%%%%%%

\begin{applicationbox}
\textbf{Numerical Analysis.}
Complicated functions and solutions are replaced by finite-dimensional
approximations suitable for computation.

\medskip

\textbf{Data Fitting.}
Least-squares approximation constructs models from noisy observations.

\medskip

\textbf{Differential Equations.}
Galerkin, spectral, and finite element methods approximate solutions of
ODEs and PDEs.

\medskip

\textbf{Signal Processing.}
Fourier and wavelet approximations represent signals using a manageable
number of coefficients.

\medskip

\textbf{Computer Graphics.}
Splines and piecewise polynomial curves are used to model smooth shapes.

\medskip

\textbf{Scientific Computing.}
Orthogonal polynomial expansions provide efficient high-order numerical
representations.

\medskip

\textbf{Optimization.}
Best approximation is itself an optimization problem over a function
class.

\medskip

\textbf{Probability and Statistics.}
Orthogonal expansions and least-squares projections appear in regression
and stochastic approximation.

\medskip

\textbf{Functional Analysis.}
Density, projection, compactness, and finite-rank approximation are
fundamental functional-analytic concepts.

\medskip

\textbf{Machine Learning.}
Function classes are studied according to their ability to approximate
target mappings from data.

\medskip

\textbf{Compression.}
Sparse Fourier and wavelet representations reduce storage while
controlling reconstruction error.
\end{applicationbox}

%%%%%%%%%%%%%%%%%%%%%%%%%%%%%%%%%%%%%%%%%%%%%%%%%%%%%%%%%%%%
\subsection{Exercises}
\label{subsec:approximation-theory-exercises}
%%%%%%%%%%%%%%%%%%%%%%%%%%%%%%%%%%%%%%%%%%%%%%%%%%%%%%%%%%%%

\begin{exercise}[1]
Explain the difference between interpolation and approximation.
\end{exercise}

\begin{exercise}[1]
Define the best approximation to \(f\) from a subset \(M\) of a normed
space.
\end{exercise}

\begin{exercise}[1]
State the Weierstrass approximation theorem.
\end{exercise}

\begin{exercise}[1]
Explain what it means for the polynomials to be dense in
\(C([a,b])\).
\end{exercise}

\begin{exercise}[1]
Write the Bernstein polynomial \(B_nf\) for a continuous function on
\([0,1]\).
\end{exercise}

\begin{exercise}[1]
Compute \(B_1f\) for

\[
f(x)=x^2.
\]
\end{exercise}

\begin{exercise}[1]
Compute \(B_2f\) for

\[
f(x)=x.
\]
\end{exercise}

\begin{exercise}[2]
Show that Bernstein operators preserve constants.
\end{exercise}

\begin{exercise}[2]
Show that

\[
B_nx=x.
\]
\end{exercise}

\begin{exercise}[2]
Compute

\[
B_n(x^2)
\]

and show that it converges uniformly to \(x^2\).
\end{exercise}

\begin{exercise}[2]
Derive the normal equations for least-squares approximation from a
finite-dimensional subspace.
\end{exercise}

\begin{exercise}[2]
Show that the orthogonal projection is the best approximation in a
Hilbert space.
\end{exercise}

\begin{exercise}[2]
Find the best constant \(L^2\) approximation to \(f(x)=x\) on
\([0,1]\).
\end{exercise}

\begin{exercise}[2]
Find the best constant uniform approximation to \(f(x)=x\) on
\([0,1]\).
\end{exercise}

\begin{exercise}[2]
Compare the answers to the previous two exercises.
\end{exercise}

\begin{exercise}[2]
Verify the recurrence

\[
T_{n+1}(x)
=
2xT_n(x)-T_{n-1}(x)
\]

for Chebyshev polynomials.
\end{exercise}

\begin{exercise}[2]
Find the zeros of \(T_3(x)\).
\end{exercise}

\begin{exercise}[2]
Show that

\[
|T_n(x)|\leq1
\]

on \([-1,1]\).
\end{exercise}

\begin{exercise}[2]
Find the best constant uniform approximation to

\[
f(x)=x^2
\]

on \([-1,1]\).
\end{exercise}

\begin{exercise}[2]
Find the least-squares constant approximation to

\[
f(x)=x^2
\]

on \([-1,1]\).
\end{exercise}

\begin{exercise}[3]
Use Gram--Schmidt to construct the first three orthogonal polynomials on
\([-1,1]\) with weight \(1\).
\end{exercise}

\begin{exercise}[3]
Derive the first three Legendre polynomials up to normalization.
\end{exercise}

\begin{exercise}[3]
Find the degree-one least-squares polynomial approximation to
\(f(x)=x^2\) on \([0,1]\).
\end{exercise}

\begin{exercise}[3]
Use the equioscillation principle to characterize the best degree-one
uniform approximation to a suitable continuous function.
\end{exercise}

\begin{exercise}[3]
Compute the modulus of continuity of

\[
f(x)=x
\]

on \([0,1]\).
\end{exercise}

\begin{exercise}[3]
Show that if \(f\) is Lipschitz with constant \(L\), then

\[
\omega(f,\delta)\leq L\delta.
\]
\end{exercise}

\begin{exercise}[3]
Compare equally spaced interpolation nodes with Chebyshev nodes for a
high-degree polynomial interpolation problem.
\end{exercise}

\begin{exercise}[3]
Construct a piecewise linear approximation to a smooth function and
estimate its error.
\end{exercise}

\begin{exercise}[3]
Explain why a cubic spline is typically required to have continuous
first and second derivatives.
\end{exercise}

\begin{exercise}[3]
Construct a Pad\'e approximant for a simple power series.
\end{exercise}

\begin{exercise}[3]
Show that mollification approximates a continuous function uniformly on
compact sets under suitable hypotheses.
\end{exercise}

\begin{exercise}[3]
Interpret a Fourier partial sum as an \(L^2\) best approximation.
\end{exercise}

\begin{exercise}[4]
Prove the Weierstrass approximation theorem using Bernstein polynomials.
\end{exercise}

\begin{exercise}[4]
Study a proof of the Stone--Weierstrass theorem.
\end{exercise}

\begin{exercise}[4]
Prove an appropriate form of the Chebyshev equioscillation theorem.
\end{exercise}

\begin{exercise}[4]
Study Jackson's theorem and derive an approximation-rate estimate from
the modulus of continuity.
\end{exercise}

\begin{exercise}[4]
Study an inverse theorem connecting polynomial approximation rates with
smoothness.
\end{exercise}

\begin{exercise}[4]
Investigate spline approximation error as the maximum knot spacing tends
to zero.
\end{exercise}

\begin{exercise}[4]
Study best \(m\)-term wavelet approximation for functions with localized
singularities.
\end{exercise}

\begin{exercise}[4]
Investigate the relationship between finite element approximation and
best approximation in Sobolev spaces.
\end{exercise}

%%%%%%%%%%%%%%%%%%%%%%%%%%%%%%%%%%%%%%%%%%%%%%%%%%%%%%%%%%%%
\subsection{Conceptual Summary}
\label{subsec:approximation-theory-summary}
%%%%%%%%%%%%%%%%%%%%%%%%%%%%%%%%%%%%%%%%%%%%%%%%%%%%%%%%%%%%

Approximation theory studies how well an object \(f\) can be represented
by simpler objects \(g\) from an approximation family.

The central quantity is

\[
\boxed{
\|f-g\|.
}
\]

A best approximation satisfies

\[
\boxed{
\|f-g^\ast\|
=
\inf_{g\in M}\|f-g\|.
}
\]

In Hilbert spaces, best approximation from a closed subspace is obtained
by orthogonal projection.

The Weierstrass theorem states that polynomials are uniformly dense in

\[
\boxed{
C([a,b]).
}
\]

Bernstein polynomials provide a constructive approximation:

\[
\boxed{
B_nf\to f.
}
\]

Stone--Weierstrass generalizes polynomial density to broad classes of
functions on compact spaces.

Least-squares approximation minimizes

\[
\boxed{
L^2\text{ error},
}
\]

while minimax approximation minimizes

\[
\boxed{
L^\infty\text{ error}.
}
\]

Chebyshev polynomials and equioscillation characterize important uniform
approximation problems.

Approximation rates depend strongly on smoothness:

\[
\boxed{
\text{more regularity}
\Longrightarrow
\text{typically faster convergence}.
}
\]

Splines replace high-degree global polynomials with stable local
piecewise polynomials.

Rational functions, Fourier expansions, wavelets, mollifiers, and
finite-dimensional projection methods provide additional approximation
frameworks.

Thus approximation theory connects the infinite-dimensional mathematical
world with finite and computational representations.

%%%%%%%%%%%%%%%%%%%%%%%%%%%%%%%%%%%%%%%%%%%%%%%%%%%%%%%%%%%%
\subsection{Section Connections}
\label{subsec:approximation-theory-connections}
%%%%%%%%%%%%%%%%%%%%%%%%%%%%%%%%%%%%%%%%%%%%%%%%%%%%%%%%%%%%

\chapterconnection{
Approximation Theory provides the bridge between exact mathematical
objects and simpler finite representations. Real Analysis supplies
uniform convergence, continuity, smoothness, and compactness. Functional
Analysis provides density, Hilbert-space projection, best approximation,
and finite-dimensional subspaces. Fourier Analysis interprets Fourier
series as orthogonal approximations and provides trigonometric
approximation. Harmonic Analysis adds localized frequency decompositions
and wavelet-type methods. Distribution Theory and Sobolev theory use
mollification to approximate rough functions by smooth functions.
Calculus of Variations and PDE theory lead to Galerkin and finite element
approximations. Optimization appears through least-squares and minimax
problems. Numerical Analysis develops stable algorithms for interpolation,
splines, spectral methods, and computational approximation. The next
section, Asymptotic Analysis, studies a related but distinct question:
how mathematical quantities can be approximated systematically in
limiting regimes.
}

%%%%%%%%%%%%%%%%%%%%%%%%%%%%%%%%%%%%%%%%%%%%%%%%%%%%%%%%%%%%
% SECTION SOURCE TRACKING
%%%%%%%%%%%%%%%%%%%%%%%%%%%%%%%%%%%%%%%%%%%%%%%%%%%%%%%%%%%%

%------------------------------------------------------------
% 5.16 Approximation Theory
%------------------------------------------------------------
%
% Source basis:
%   - Standard classical approximation theory.
%   - Standard functional analytic best-approximation theory.
%   - Standard polynomial, spline, Fourier, and rational approximation.
%
% Used for:
%   - approximation errors and norms;
%   - interpolation versus approximation;
%   - best approximation;
%   - orthogonal projection;
%   - least-squares approximation;
%   - polynomial approximation;
%   - Weierstrass approximation theorem;
%   - Bernstein polynomials;
%   - positive linear operators;
%   - Korovkin-type approximation;
%   - Stone--Weierstrass theorem;
%   - trigonometric approximation;
%   - Fourier approximation;
%   - orthogonal polynomials;
%   - Legendre and Chebyshev polynomials;
%   - Chebyshev nodes;
%   - Runge phenomenon;
%   - minimax approximation;
%   - equioscillation;
%   - approximation rates;
%   - modulus of continuity;
%   - Jackson-type estimates;
%   - inverse approximation theorems;
%   - piecewise polynomial approximation;
%   - splines and B-splines;
%   - rational and Pade approximation;
%   - convolution and mollification;
%   - Sobolev approximation;
%   - Galerkin approximation;
%   - finite element connections;
%   - nonlinear and sparse approximation;
%   - wavelet approximation;
%   - stability and conditioning.
%
% Notes:
%   - Numerical and Computational Mathematics develops interpolation,
%     numerical approximation, finite elements, and implementation in
%     greater computational detail.
%   - Fourier Analysis develops trigonometric approximation and
%     orthogonal expansions.
%   - Functional Analysis supplies the abstract projection and density
%     framework.
%   - Distribution Theory supplies mollification and smooth
%     approximation of generalized functions.
%   - Asymptotic Analysis next studies approximations arising from
%     limiting parameters and scale separation.
%
%%%%%%%%%%%%%%%%%%%%%%%%%%%%%%%%%%%%%%%%%%%%%%%%%%%%%%%%%%%%